% 블루프린트 본문. 웹 판(web.tex)과 PDF 판(print.tex)이 함께 쓴다.
% 관례:
%   \lean{...}    대응하는 Lean 선언 (namespace 포함)
%   \leanok       진술(환경 안) 또는 증명(proof 안)이 Lean 에서 검증됨
%   \uses{...}    의존하는 항목. 진술 안이면 진술에, proof 안이면 증명에만 필요
%   \notready     블루프린트를 더 다듬어야 함 (증명 계획이 아직 없음)
% 아직 증명하지 않은 가정 (H3, IT(q) 등) 은 Lean 에서 `def ... : Prop` 으로 정의만 하고,
% 조건부 정리의 명시적 가정으로 쓴다. 문헌 정리는 \lean 없이 별도 노드로 둔다.
% 각 항목 끝의 [B1], [C4] 같은 표시는 docs/lemma-ledger.md 의 ID 다.

\chapter{개요}

처음 \(n\) 개 소수를 이어 붙인 수를 \(W_n\) 이라 하자.
\(W_1 = 2,\ W_2 = 23,\ W_3 = 235,\ W_4 = 2357,\ W_5 = 235711, \dots\)
목표는 \(W_n\) 이 소수인 \(n\) 이 무한히 많다는 추측 (정리~\ref{thm:main}) 이다.
이 추측은 미해결이며, 알려진 소수 항의 첨자는 \(1, 2, 4, 128, 174, 342, 435, 1429\) 뿐이다.

이 문서는 추측으로 가는 길에서 지금까지 Lean 으로 검증한 것, 문헌에 기대는 것,
아직 계획이 없는 것을 한 그래프에 모은다. 표기는 다음과 같다.
\begin{itemize}
  \item \(p_k\) 는 \(k\) 번째 소수 (\(p_1 = 2\)). Lean 에서 \(p_{n+1}\) 은 \texttt{Nat.nth Nat.Prime n} 이다.
  \item \(\dig(m)\) 은 \(m\) 의 십진 자릿수.
  \item ``Lean 검증'' 은 \texttt{sorry} 와 추가 공리 없이 컴파일되고 \texttt{\#print axioms} 가
        \texttt{propext}, \texttt{Classical.choice}, \texttt{Quot.sound} 의 부분집합만 보이는 것을 뜻한다
        (\texttt{SmarandacheWellin/Audit.lean}).
\end{itemize}

현재 위치를 한 문장으로 말하면 이렇다. 유한 소수 집합 \(S\) 가 어느 시점 이후의 모든 \(W_n\) 을
나누어 버릴 수 있는가 (Q1) 를 묻고 있고, \(|S| = 1\) 은 무조건, \(|S| = 2\) 는 연속 소수에 관한
문헌의 정리 (Banks--Freiberg--Turnage-Butterbaugh 2015, Maynard 2016) 를 가정하면 해결했다.
\(|S| \ge 3\) 과 추측 자체는 열려 있다.

\chapter{정의와 기본 성질}

\section{정의}

\begin{definition}[자릿수]
  \label{def:digitsLen}
  \lean{SW.digitsLen}
  \leanok
  \(\dig(m)\) 은 \(m\) 의 십진 자릿수 목록의 길이다. 특히 \(\dig(0) = 0\).
\end{definition}

\begin{definition}[Smarandache--Wellin 수]
  \label{def:W}
  \lean{SW.W}
  \leanok
  \uses{def:digitsLen}
  \(W_0 = 0\), \(W_{n+1} = 10^{\dig(p_{n+1})} W_n + p_{n+1}\).
\end{definition}

\begin{definition}[목표 명제]
  \label{def:SWInfinitude}
  \lean{SW.SWInfinitude}
  \leanok
  \uses{def:W}
  모든 \(N\) 에 대해 \(N < n\) 이고 \(W_n\) 이 소수인 \(n\) 이 있다.
\end{definition}

\section{기본 성질}

\begin{lemma}
  \label{lem:basic-primes}
  \lean{SW.nth_prime_prime, SW.digitsLen_pos}
  \leanok
  \uses{def:digitsLen}
  \(p_{n+1}\) 은 소수이고, \(m \ne 0\) 이면 \(\dig(m) > 0\).
\end{lemma}
\begin{proof}
  \leanok
  Mathlib 의 \texttt{Nat.nth\_mem\_of\_infinite} 와 자릿수 목록이 비어 있지 않다는 사실.
\end{proof}

\begin{lemma}[점화식]
  \label{lem:W-rec}
  \lean{SW.W_zero, SW.W_succ}
  \leanok
  \uses{def:W}
  \(W_0 = 0\), \(W_{n+1} = 10^{\dig(p_{n+1})} W_n + p_{n+1}\).
\end{lemma}
\begin{proof}
  \leanok
  정의 그대로.
\end{proof}

\begin{lemma}[처음 네 항, B1]
  \label{lem:W-small}
  \lean{SW.W_one, SW.W_two, SW.W_three, SW.W_four, SW.W_one_prime, SW.W_two_prime, SW.W_three_not_prime, SW.W_four_prime}
  \leanok
  \uses{def:W}
  \(W_1 = 2\), \(W_2 = 23\), \(W_3 = 235\), \(W_4 = 2357\). 이 가운데 \(2, 23, 2357\) 은 소수이고 \(235\) 는 합성수다.
  정의가 OEIS A019518, A046035 와 맞는다는 확인이다.
\end{lemma}
\begin{proof}
  \leanok
  \(p_1, \dots, p_4\) 를 Mathlib 의 값으로 바꿔 계산한다.
\end{proof}

\begin{lemma}[증가, B2]
  \label{lem:W-mono}
  \lean{SW.W_lt_succ, SW.W_strictMono, SW.W_injective, SW.W_pos}
  \leanok
  \uses{lem:W-rec, lem:basic-primes}
  \(W_n < W_{n+1}\). 따라서 \(W\) 는 순증가, 단사이고 \(n \ge 1\) 이면 \(W_n > 0\).
\end{lemma}
\begin{proof}
  \leanok
  \(10^{\dig(p_{n+1})} \ge 1\) 이고 \(p_{n+1} > 0\).
\end{proof}

\begin{lemma}[나머지 점화식, B3, B4]
  \label{lem:W-mod}
  \lean{SW.W_succ_mod, SW.W_succ_modEq_ten}
  \leanok
  \uses{lem:W-rec, lem:basic-primes}
  \(W_{n+1} \bmod m = \bigl((10^{\dig(p_{n+1})} \bmod m)(W_n \bmod m) + p_{n+1} \bmod m\bigr) \bmod m\).
  특히 \(W_{n+1} \equiv p_{n+1} \pmod{10}\).
\end{lemma}
\begin{proof}
  \leanok
  나머지 산술. 두 번째는 \(10 \mid 10^{\dig(p_{n+1})}\) 에서.
\end{proof}

\begin{theorem}[목표 명제의 다른 서술, B5]
  \label{thm:infinitude-iff}
  \lean{SW.swInfinitude_iff_infinite, SW.swInfinitude_iff_image}
  \leanok
  \uses{def:SWInfinitude}
  목표 명제는 \(\{n : W_n \text{ 소수}\}\) 가 무한집합이라는 것, 또 서로 다른 Smarandache--Wellin 소수가
  무한히 많다는 것과 같다.
\end{theorem}
\begin{proof}
  \leanok
  \uses{lem:W-mono}
  첫째는 \(\exists^\infty\) 의 표준 변환. 둘째는 \(W\) 가 단사이므로.
\end{proof}

\section{작은 모듈러스}

\begin{theorem}[\(10 \equiv 1\) 인 모듈러스, C1]
  \label{thm:mod-sum}
  \lean{SW.W_modEq_sum_of_ten_modEq_one, SW.W_modEq_sum_three, SW.W_modEq_sum_nine}
  \leanok
  \uses{def:W}
  \(10 \equiv 1 \pmod m\) 이면 \(W_n \equiv p_1 + \dots + p_n \pmod m\). 특히 \(m = 3, 9\).
\end{theorem}
\begin{proof}
  \leanok
  \uses{lem:W-rec}
  \(n\) 에 대한 귀납법. \(10^{\dig} \equiv 1\) 이므로 이어 붙이기가 더하기가 된다.
\end{proof}

\begin{lemma}[한 소수는 이웃한 두 항을 못 나눈다, C2]
  \label{lem:not-dvd-succ}
  \lean{SW.not_dvd_succ_of_dvd, SW.not_dvd_succ_of_dvd_of_lt}
  \leanok
  \uses{def:W}
  소수 \(q\) 가 \(W_n\) 을 나누고 \(q \ne p_{n+1}\) 이면 \(q \nmid W_{n+1}\).
\end{lemma}
\begin{proof}
  \leanok
  \uses{lem:W-rec, lem:basic-primes}
  \(q\) 가 \(W_n\) 과 \(W_{n+1} = 10^{\dig} W_n + p_{n+1}\) 을 모두 나누면 \(q \mid p_{n+1}\) 이고,
  두 소수이므로 \(q = p_{n+1}\).
\end{proof}

\begin{lemma}[\(2, 5\) 의 자동 배제, C3]
  \label{lem:coprime-ten}
  \lean{SW.seven_le_nth_prime, SW.nth_prime_coprime_ten, SW.coprime_ten_of_three_le}
  \leanok
  \uses{def:W}
  \(n \ge 3\) 이면 \(p_{n+1} \ge 7\) 이고 \(\gcd(W_{n+1}, 10) = 1\).
\end{lemma}
\begin{proof}
  \leanok
  \uses{lem:W-mod, lem:basic-primes}
  끝자리가 \(p_{n+1}\) 의 끝자리이고, \(p_{n+1}\) 은 \(2, 5\) 가 아닌 소수다.
\end{proof}

\chapter{유한 소수 집합 회피 (Q1)}

질문 Q1: 유한 소수 집합 \(S\) 에 대해, \(S\) 의 어느 원소도 \(W_n\) 을 나누지 않는 \(n\) 이 무한히 많은가?
그렇지 않다면 어느 시점 이후 모든 \(W_n\) 이 \(S\) 의 어떤 원소로 나누어지고 (\(S\) 가 ``덮는다''),
목표 추측은 거짓이다. 그래서 Q1 은 추측의 필요조건이다.

\section{두 칸 점화식}

\begin{lemma}[첨자와 소수의 크기]
  \label{lem:lt-nth-prime}
  \lean{SW.lt_nth_prime, SW.not_dvd_nth_prime_of_le, SW.lt_of_nth_prime_lt}
  \leanok
  \uses{lem:basic-primes}
  \(n < p_{n+1}\). 따라서 소수 \(q \le n\) 은 \(p_{n+1}\) 을 나누지 않는다.
  또 \(p_{M+1} < p_{m+1}\) 이면 \(M < m\).
\end{lemma}
\begin{proof}
  \leanok
  \(p\) 가 순증가하므로 귀납법.
\end{proof}

\begin{definition}[블록]
  \label{def:B}
  \lean{SW.B}
  \leanok
  \uses{def:digitsLen}
  \(B_n = 10^{\dig(p_{n+2})} p_{n+1} + p_{n+2}\). 연속한 두 소수를 이어 붙인 수다. 예: \(B_6 = 1719\).
\end{definition}

\begin{lemma}[두 칸 점화식, C4]
  \label{lem:W-add-two}
  \lean{SW.W_add_two, SW.dvd_add_two_iff}
  \leanok
  \uses{def:W, def:B}
  \(W_{n+2} = 10^{\dig(p_{n+2}) + \dig(p_{n+1})} W_n + B_n\).
  따라서 \(q \mid W_n\) 이면 \(q \mid W_{n+2} \iff q \mid B_n\).
\end{lemma}
\begin{proof}
  \leanok
  \uses{lem:W-rec}
  점화식을 두 번 쓴다.
\end{proof}

\section{소수 하나}

\begin{theorem}[Q1, \(|S| = 1\), C5]
  \label{thm:one-prime}
  \lean{SW.not_dvd_of_three_consecutive, SW.exists_not_dvd}
  \leanok
  \uses{def:W}
  소수 \(q\) 는 연속한 세 항 \(W_{n+1}, W_{n+2}, W_{n+3}\) 을 모두 나누지는 못한다.
  따라서 \(q \nmid W_n\) 인 \(n\) 이 무한히 많다.
\end{theorem}
\begin{proof}
  \leanok
  \uses{lem:not-dvd-succ}
  셋 다 나누면 보조정리~\ref{lem:not-dvd-succ} 에 의해 \(q = p_{n+2} = p_{n+3}\) 이 되어 모순.
\end{proof}

\section{소수 둘: 덮개는 번갈아 나눈다}

\begin{lemma}[덮개와 블록]
  \label{lem:cover-block}
  \lean{SW.block_dvd_of_cover, SW.block_dvd_or_succ_of_cover}
  \leanok
  \uses{def:W, def:B}
  소수 \(q, r\) 가 \(N\) 이후의 모든 항을 덮는다고 하자. 즉 \(n > N\) 이면 \(q \mid W_n\) 또는 \(r \mid W_n\).
  그러면 \(q, r \le n\) 인 모든 \(n > N\) 에서 \(q \mid B_n\) 또는 \(q \mid B_{n+1}\) 이다.
\end{lemma}
\begin{proof}
  \leanok
  \uses{lem:not-dvd-succ, lem:W-add-two, lem:lt-nth-prime}
  보조정리~\ref{lem:not-dvd-succ} 에 의해 한 소수는 이웃한 두 항을 못 나누므로, 덮으려면 \(q, r\) 가
  번갈아 나눠야 한다. 그러면 \(q\) 는 \(n\) 또는 \(n+1\) 에서 두 칸 간격으로 나누고,
  보조정리~\ref{lem:W-add-two} 로 그 사이 블록을 나눈다.
\end{proof}

\begin{theorem}[블록 조건에서 Q1, \(|S| = 2\), C6]
  \label{thm:two-of-block}
  \lean{SW.avoid_two_of_block, SW.avoid_two_of_block_succ}
  \leanok
  \uses{def:W, def:B}
  소수 \(q\) 에 대해 \(q \nmid B_n\) 이고 \(q \nmid B_{n+1}\) 인 \(n\) 이 무한히 많으면 (가정 HB(\(q\))),
  임의의 소수 \(r\) 에 대해 \(q \nmid W_n\) 이고 \(r \nmid W_n\) 인 \(n\) 이 무한히 많다.
\end{theorem}
\begin{proof}
  \leanok
  \uses{lem:cover-block}
  덮개가 있다면 보조정리~\ref{lem:cover-block} 와 모순.
\end{proof}

\section{\(q = 3\): 같은 류의 연속 소수 세 개}

\begin{lemma}[\(3\) 과 블록, C7]
  \label{lem:three-B}
  \lean{SW.three_dvd_B_iff, SW.not_three_dvd_pow_add_one, SW.not_three_dvd_B_of_modEq}
  \leanok
  \uses{def:B}
  \(3 \mid B_n \iff 3 \mid p_{n+1} + p_{n+2}\). 또 \(3 \nmid 10^d + 1\).
  따라서 \(n \ge 3\) 이고 \(p_{n+1} \equiv p_{n+2} \pmod 3\) 이면 \(3 \nmid B_n\).
\end{lemma}
\begin{proof}
  \leanok
  \uses{lem:lt-nth-prime}
  \(10 \equiv 1 \pmod 3\). 같은 류이면 합은 \(2 p_{n+1}\) 이고 \(3 \nmid p_{n+1}\).
\end{proof}

\begin{definition}[같은 류의 세 쌍]
  \label{def:same-class-triple}
  \lean{SW.SameClassTriple}
  \leanok
  \(p_{m+1} \equiv p_{m+2} \equiv p_{m+3} \pmod q\).
\end{definition}

\begin{definition}[가정 H3]
  \label{def:H3}
  \lean{SW.SameClassTriples3}
  \leanok
  \uses{def:same-class-triple}
  모든 \(M\) 에 대해 \(M < m\) 이고 \(p_{m+1} \equiv p_{m+2} \equiv p_{m+3} \pmod 3\) 인 \(m\) 이 있다.
  \(W\) 와 무관한, 소수만의 명제다.
\end{definition}

\begin{theorem}[H3 에서 Q1, \(S = \{3, r\}\), C7]
  \label{thm:three-of-triples}
  \lean{SW.avoid_three_of_triples}
  \leanok
  \uses{def:H3, def:W}
  H3 이 성립하면 임의의 소수 \(r\) 에 대해 \(3 \nmid W_n\) 이고 \(r \nmid W_n\) 인 \(n\) 이 무한히 많다.
\end{theorem}
\begin{proof}
  \leanok
  \uses{lem:three-B, thm:two-of-block}
  같은 류의 세 쌍은 \(3 \nmid B_m\), \(3 \nmid B_{m+1}\) 을 주므로 HB(3) 이 성립한다.
\end{proof}

\section{일반 \(q\): 자릿수 두 구간}

\begin{lemma}[같은 류 쌍의 블록, C8]
  \label{lem:pow-add-one}
  \lean{SW.dvd_pow_add_one_of_dvd_B, SW.not_dvd_pow_add_one_both}
  \leanok
  \uses{def:B}
  \(p_{n+1} \equiv p_{n+2} \pmod q\), \(q \nmid p_{n+1}\), \(q \mid B_n\) 이면 \(q \mid 10^{\dig(p_{n+2})} + 1\).
  또 어떤 소수도 \(10^d + 1\) 과 \(10^{d+1} + 1\) 을 동시에 나누지 못한다.
\end{lemma}
\begin{proof}
  \leanok
  \uses{lem:three-B}
  \(B_n \equiv p_{n+1}(10^d + 1) \pmod q\). 둘을 다 나누면 \(q \mid 10(10^d+1) - (10^{d+1}+1) = 9\) 이므로
  \(q = 3\) 인데 \(3 \nmid 10^d + 1\).
\end{proof}

\begin{lemma}[덮개 아래의 자릿수 조건]
  \label{lem:cover-pow}
  \lean{SW.dvd_pow_add_one_of_cover}
  \leanok
  \uses{def:same-class-triple, def:W}
  \(q, r\) 가 \(N\) 이후를 덮고, \(m > N\), \(q, r \le m\) 에서 mod \(q\) 같은 류의 세 쌍이 시작하며
  \(p_{m+2}, p_{m+3}\) 가 모두 \(d\) 자리이면 \(q \mid 10^d + 1\).
\end{lemma}
\begin{proof}
  \leanok
  \uses{lem:cover-block, lem:pow-add-one, lem:lt-nth-prime}
  보조정리~\ref{lem:cover-block} 로 \(q \mid B_m\) 또는 \(q \mid B_{m+1}\), 그 뒤 보조정리~\ref{lem:pow-add-one}.
\end{proof}

\begin{definition}[가정 Hd(\(q\))]
  \label{def:Hd}
  \lean{SW.DigitTriples}
  \leanok
  \uses{def:same-class-triple, def:digitsLen}
  모든 \(M\) 에 대해 어떤 \(d\) 가 있어, \(M\) 뒤에 mod \(q\) 같은 류의 세 쌍 가운데
  \(p_{m+2}, p_{m+3}\) 가 \(d\) 자리인 것과 \(d+1\) 자리인 것이 각각 있다.
\end{definition}

\begin{theorem}[Hd(\(q\)) 에서 Q1, \(S = \{q, r\}\), C9]
  \label{thm:two-of-digit}
  \lean{SW.avoid_two_of_digit_triples}
  \leanok
  \uses{def:Hd, def:W}
  Hd(\(q\)) 가 성립하면 임의의 소수 \(r\) 에 대해 \(q \nmid W_n\) 이고 \(r \nmid W_n\) 인 \(n\) 이 무한히 많다.
\end{theorem}
\begin{proof}
  \leanok
  \uses{lem:cover-pow, lem:pow-add-one}
  덮개가 있다면 \(q \mid 10^d + 1\) 과 \(q \mid 10^{d+1} + 1\) 이 함께 나와 모순.
\end{proof}

\section{구간 판 가정 IT(\(q\))}

\begin{lemma}[자릿수 구간]
  \label{lem:digitsLen-eq}
  \lean{SW.digitsLen_eq}
  \leanok
  \uses{def:digitsLen}
  \(10^e \le n < 10^{e+1}\) 이면 \(\dig(n) = e + 1\).
\end{lemma}
\begin{proof}
  \leanok
  Mathlib 의 \texttt{Nat.digits\_length\_le\_iff}, \texttt{Nat.lt\_digits\_length\_iff}.
\end{proof}

\begin{definition}[가정 IT(\(q\))]
  \label{def:IT}
  \lean{SW.IntervalTriples}
  \leanok
  \uses{def:same-class-triple}
  어떤 \(x_0\) 가 있어 \(x \ge x_0\) 이면 \(x \le p_{m+1}\), \(p_{m+3} \le 3x\) 이고
  \(p_{m+1} \equiv p_{m+2} \equiv p_{m+3} \pmod q\) 인 \(m\) 이 있다.
\end{definition}

\begin{lemma}[IT(\(q\)) 에서 Hd(\(q\)), H3]
  \label{lem:IT-Hd}
  \lean{SW.digit_triple_of_interval, SW.digitTriples_of_intervalTriples, SW.sameClassTriples3_of_intervalTriples}
  \leanok
  \uses{def:IT, def:Hd, def:H3}
  IT(\(q\)) 이면 Hd(\(q\)) 이다. IT(3) 이면 H3 이다.
\end{lemma}
\begin{proof}
  \leanok
  \uses{lem:digitsLen-eq, lem:lt-nth-prime}
  \([10^e, 3 \cdot 10^e] \subset [10^e, 10^{e+1})\) 의 수는 모두 \(e + 1\) 자리다. \(x = 10^e, 10^{e+1}\) 에 적용한다.
\end{proof}

\begin{theorem}[IT(\(q\)) 에서 Q1, \(S = \{q, r\}\)]
  \label{thm:two-of-IT}
  \lean{SW.avoid_two_of_interval_triples, SW.avoid_three_of_interval_triples}
  \leanok
  \uses{def:IT, def:W}
  IT(\(q\)) 가 성립하면 임의의 소수 \(r\) 에 대해 \(q \nmid W_n\) 이고 \(r \nmid W_n\) 인 \(n\) 이 무한히 많다.
\end{theorem}
\begin{proof}
  \leanok
  \uses{lem:IT-Hd, thm:two-of-digit, thm:three-of-triples}
  보조정리~\ref{lem:IT-Hd} 와 정리~\ref{thm:two-of-digit}.
\end{proof}

\section{문헌의 정리 (Lean 밖)}

다음 두 정리는 Maynard 체에 기반한 깊은 결과이며 Mathlib 에 없다. 이 프로젝트는 이들을 형식화하지 않고,
위의 조건부 정리의 가정으로만 쓴다. 대조 기록은 \texttt{docs/literature/2026-11-05-shiu-bft-maynard.md}.

\begin{theorem}[문헌: Banks--Freiberg--Turnage-Butterbaugh 2015, Corollary 3]
  \label{thm:lit-BFT}
  \uses{def:H3}
  H3 이 성립한다. 원문은 더 강하다: \((a, D) = 1\), \(D \ge 3\), \(m \ge 2\) 이면
  \(p_{r+1} \equiv \dots \equiv p_{r+m} \equiv a \pmod D\) 인 \(r\) 이 무한히 많다. \(D = 3\), \(m = 3\) 이 H3 이다.
\end{theorem}
\begin{proof}
  W.~D.~Banks, T.~Freiberg, C.~L.~Turnage-Butterbaugh, \emph{Consecutive primes in tuples},
  Acta Arith.~167 (2015), 261--266, arXiv:1311.7003. Shiu (2000) 의 임의 길이 문자열에서도 따른다.
\end{proof}

\begin{theorem}[문헌: Maynard 2016, Theorem 3.3 의 증명 중 식 (6.24)]
  \label{thm:lit-Maynard}
  \uses{def:IT}
  모든 소수 \(q \ge 3\) 에 대해 IT(\(q\)) 가 성립한다.
\end{theorem}
\begin{proof}
  J.~Maynard, \emph{Dense clusters of primes in subsets}, Compositio Math.~152 (2016), 1517--1554.
  식 (6.24): \(x > x_0(\varepsilon, q)\) 인 모든 \(x\) 에 대해 \(x \le n \le 2x\) 이고
  \([qn, qn + \eta \log x]\) 안에 mod \(q\) 같은 류의 연속 소수가 \(\gg \log k\) 개 있는 \(n\) 의 수가
  \(\gg x / (q^k \exp(C_6 k^5))\). \(x\) 를 \(x/q\) 로 바꾸면 IT(\(q\)) 다.
  주의: 정리문 (Theorem 3.3) 자체는 \([x, Tx]\) (\(T\) 는 상수) 형태까지만 주며, 이 노드는 증명 내부 식을 인용한다.
\end{proof}

\begin{corollary}[Q1, \(|S| = 2\), 문헌 조건부]
  \label{cor:Q1-two}
  \uses{def:W}
  모든 소수 \(q, r\) 에 대해 \(q \nmid W_n\) 이고 \(r \nmid W_n\) 인 \(n\) 이 무한히 많다.
\end{corollary}
\begin{proof}
  \uses{thm:lit-Maynard, thm:two-of-IT, lem:coprime-ten, thm:one-prime}
  \(q \ge 3\) 이고 \(q \ne 5\) 이면 정리~\ref{thm:lit-Maynard} 와 정리~\ref{thm:two-of-IT}.
  \(q \in \{2, 5\}\) 이면 \(n \ge 4\) 에서 \(q \nmid W_n\) (보조정리~\ref{lem:coprime-ten}) 이므로
  \(r\) 하나의 경우 (정리~\ref{thm:one-prime}) 로 줄어든다. \(r\) 에 대해서도 대칭.
  Lean 에서는 문헌 정리를 가정으로 두는 형태 (정리~\ref{thm:two-of-IT}) 까지만 검증했다.
\end{proof}

\section{열린 문제}

\begin{proposition}[Q1, \(|S| \ge 3\)]
  \label{prop:Q1-general}
  \notready
  \uses{def:W}
  임의의 유한 소수 집합 \(S\) 에 대해 \(S\) 의 어느 원소도 \(W_n\) 을 나누지 않는 \(n\) 이 무한히 많다.
\end{proposition}
\begin{proof}
  \uses{thm:one-prime, cor:Q1-two}
  계획 없음. 세 소수 이상이 덮을 때는 번갈아 나눌 필요가 없어 두 칸 구조가 사라진다.
  \(3 \in S\) 이면 \(W_n \bmod 3\) 이 소수 합 \(\sum p_k \bmod 3\) 의 분포 문제가 된다.
  계산상으로는 \(n \le 2 \cdot 10^5\) 에서 회피 비율이 독립 모형 \(\prod_{q \in S}(1 - 1/q)\) 과 0.3\% 이내로 맞는다.
\end{proof}

\begin{proposition}[Q2, 나머지의 균등분포]
  \label{prop:Q2}
  \notready
  \uses{def:W}
  \(\gcd(m, 10) = 1\) 인 고정된 \(m\) 에 대해 \(W_n \bmod m\) 의 방문 빈도는 균등하다.
\end{proposition}
\begin{proof}
  \uses{lem:W-mod}
  계획 없음. 각 단계가 \(\mathbb{Z}/m\) 의 순열이어도 궤도가 균등분포한다는 보장은 없다.
\end{proof}

\begin{proposition}[Q4, 큰 소인수의 곱 배제]
  \label{prop:Q4}
  \notready
  \uses{def:W}
  분포 정보로 ``남은 \(W_n\) 이 큰 소수들의 곱일 가능성'' 을 배제한다.
\end{proposition}
\begin{proof}
  \uses{prop:Q1-general, prop:Q2}
  계획 없음. 일반 체 방법의 패리티 장벽이 가로막는다.
\end{proof}

\begin{theorem}[목표 추측, G]
  \label{thm:main}
  \notready
  \uses{def:SWInfinitude}
  \(W_n\) 이 소수인 \(n\) 이 무한히 많다.
\end{theorem}
\begin{proof}
  \uses{prop:Q4, thm:infinitude-iff}
  미해결. 확률 모형 (자릿수 \(D_n\) 의 수가 소수일 확률 \(\approx 2.5 / (D_n \ln 10)\)) 의 기대 개수는
  \(\log \log N\) 처럼 발산하지만, 다음 소수 항은 \(n \approx 10^{10}\)--\(10^{12}\) 에서나 기대되어 계산 탐색은 불가능하다.
\end{proof}

\chapter{계산 기록}

계산은 증명이 아니다. 각 기록은 날짜, 명령, 해석을 담고 \texttt{docs/experiments/} 에 있다.
\begin{itemize}
  \item \texttt{2026-11-04-swmod-200k.txt}: \(W_n \bmod m\) 의 방문 빈도 (Q2).
  \item \texttt{2026-11-05-avoid-200k.txt}: 유한 집합 회피 비율과 블록 조건 비율 (Q1).
  \item \texttt{2026-11-05-heuristic-4e6.txt}: 확률 모형의 보정과 외삽 (G).
  \item \texttt{2026-11-05-triples-2M.txt}: H3, HB(\(q\)), Hd(\(q\)) 의 비율 (\(n < 2 \cdot 10^6\)).
\end{itemize}
